% Source: physical PDF p. 225 (printed p. 202). Coverage: complete Figure
% 28.1 and caption. Uncertainty: none; topology, particle labels, line
% styles, and arrow directions were reconstructed from the full-resolution
% rendered scan.
\begin{figure}[H]
  \centering
  \begin{tikzpicture}[
    x=1cm,
    y=1cm,
    line cap=round,
    line join=round,
    every node/.style={inner sep=1pt},
    fermion/.style={
      draw,
      line width=0.78pt,
      postaction={decorate},
      decoration={markings,
        mark=at position 0.50 with
          {\arrow{Stealth[length=2.2mm,width=1.6mm]}}}
    },
    antifermion/.style={
      draw,
      line width=0.78pt,
      postaction={decorate},
      decoration={markings,
        mark=at position 0.50 with
          {\arrowreversed{Stealth[length=2.2mm,width=1.6mm]}}}
    },
    scalar/.style={
      draw,
      line width=0.78pt,
      dash pattern=on 3.4pt off 3.4pt,
      postaction={decorate},
      decoration={markings,
        mark=at position 0.50 with
          {\arrow{Stealth[length=2.2mm,width=1.6mm]}}}
    },
    antiscalar/.style={
      draw,
      line width=0.78pt,
      dash pattern=on 3.4pt off 3.4pt,
      postaction={decorate},
      decoration={markings,
        mark=at position 0.50 with
          {\arrowreversed{Stealth[length=2.2mm,width=1.6mm]}}}
    },
    gaugino/.style={
      draw,
      line width=0.72pt,
      decorate,
      decoration={snake,amplitude=0.105cm,segment length=0.31cm}
    }
  ]
    % External quarks and the two horizontal squark propagators.
    \draw[fermion] (-5.1,1.15) -- (-2.15,1.15);
    \draw[scalar] (-2.15,1.15) -- (2.15,1.15);
    \draw[fermion] (2.15,1.15) -- (5.1,1.15);
    \draw[antifermion] (-5.1,-1.15) -- (-2.15,-1.15);
    \draw[antiscalar] (-2.15,-1.15) -- (2.15,-1.15);
    \draw[antifermion] (2.15,-1.15) -- (5.1,-1.15);

    % The two vertical propagators are combined solid and wavy gluino lines.
    \draw[line width=0.78pt] (-2.15,1.15) -- (-2.15,-1.15);
    \draw[gaugino] (-2.15,1.15) -- (-2.15,-1.15);
    \draw[line width=0.78pt] (2.15,1.15) -- (2.15,-1.15);
    \draw[gaugino] (2.15,1.15) -- (2.15,-1.15);

    \node[above] at (-4.35,1.20) {$d_L$};
    \node[above] at (4.35,1.20) {$s_L$};
    \node[below] at (-4.35,-1.20) {$s_L$};
    \node[below] at (4.35,-1.20) {$d_L$};

    \node[above] at (-1.42,1.18) {$\mathcal{D}_1$};
    \node[above] at (0,1.18) {$\mathcal{D}_2$};
    \node[above] at (1.42,1.18) {$\mathcal{D}_3$};
    \node[below] at (-1.42,-1.18) {$\mathcal{D}_1$};
    \node[below] at (0,-1.18) {$\mathcal{D}_2$};
    \node[below] at (1.42,-1.18) {$\mathcal{D}_3$};

    \node[anchor=east] at (-2.62,0) {gluino};
    \node[anchor=west] at (2.62,0) {gluino};
  \end{tikzpicture}
  \renewcommand{\thefigure}{28.1}
  \caption{A one-loop diagram that can contribute to the
  \(\Delta S=2\) effective interaction
  \((\bar{s}_L\gamma^\mu d_L)(\bar{s}_L\gamma_\mu d_L)\) in the
  supersymmetric standard model. Here solid lines are quarks; dashed
  lines are squarks; and combined solid and wavy lines are gluinos.}
  \label{fig:28.1}
\end{figure}
